% GENERATED FILE. Edit canonical lessons, then run npm run generate:books.
% canonical-source-hash: fnv1a64:664d9ffdd7b3aa0f
% canonical-lessons: ML-C335-birudam, ML-C335-visayam, ML-C335-kanakku, ML-C335-grhapatham, ML-C335-aksaram

\chapter{Degree, Subject, Mathematics, Homework, Letter}
\label{ch:ml-degree-subject-mathematics-homework-letter}
\hlchaptermodality{335}

\begin{chapteropening}
\textbf{By the end of this chapter:} \emph{I can say a degree, a subject, mathematics, homework and a letter.}

Complete the last lesson of chapter 335: I can say a degree, a subject, mathematics, homework and a letter.
\end{chapteropening}

\section[\texorpdfstring{\ml{ബിരുദം} (birudaṁ)}{ബിരുദം (birudaṁ)}]{\ml{ബിരുദം} (birudaṁ) — a degree}

\label{lesson:ML-C335-birudam}



\begin{quote}
Before the new one: say the Malayalam for a company, then the Malayalam for a meeting.
\end{quote}

\subsection*{You'll want to know: \ml{ബിരുദം}}
\textbf{\ml{ബിരുദം}} — \emph{birudaṁ} — \textquotedblleft{}a degree\textquotedblright{}.

\begin{cousinweb}
\textbf{\ml{ബിരുദം}} (\emph{birudaṁ}) comes from Sanskrit \emph{biruda}, \textquotedblleft{}a title, a mark of honour\textquotedblright{}.
\end{cousinweb}

\begin{grammarlens}[title={What you've built}]
One more word for this chapter.
\end{grammarlens}

\subsection*{Your turn}
\begin{itemize}
\raggedright
  \item \emph{Say it:} \emph{birudaṁ}
  \item \emph{Say it:} \emph{birudaṁ}, once more
  \item \emph{Say it:} say \emph{birudaṁ}
  \item \emph{Recall:} say the Malayalam for a company, then the Malayalam for a meeting, then say \emph{birudaṁ} again
\end{itemize}

\begin{tcolorbox}[breakable,colback=teal!4,colframe=teal!35!black,title={Before you move on}]
What does \ml{ബിരുദം} mean? (\textquotedblleft{}A degree\textquotedblright{}.) Say it once more.
\end{tcolorbox}
\section[\texorpdfstring{\ml{വിഷയം} (viṣayaṁ)}{വിഷയം (viṣayaṁ)}]{\ml{വിഷയം} (viṣayaṁ) — a subject, a topic}

\label{lesson:ML-C335-visayam}



\begin{quote}
Before the new one: say the Malayalam for a meeting, then the Malayalam for a degree.
\end{quote}

\subsection*{You'll want to know: \ml{വിഷയം}}
\textbf{\ml{വിഷയം}} — \emph{viṣayaṁ} — \textquotedblleft{}a subject, a topic\textquotedblright{}.

\begin{cousinweb}
\textbf{\ml{വിഷയം}} (\emph{viṣayaṁ}) $\leftarrow$ Sanskrit \emph{viṣaya}, \textquotedblleft{}scope, sphere\textquotedblright{}, taken over whole as a tatsama, with Malayalam's noun ending \emph{-ṁ} (CDIAL 11973).
\end{cousinweb}

\begin{grammarlens}[title={What you've built}]
One more word for this chapter.
\end{grammarlens}

\subsection*{Your turn}
\begin{itemize}
\raggedright
  \item \emph{Say it:} \emph{viṣayaṁ}
  \item \emph{Say it:} \emph{viṣayaṁ}, once more
  \item \emph{Say it:} say \emph{viṣayaṁ}
  \item \emph{Recall:} say the Malayalam for a meeting, then the Malayalam for a degree, then say \emph{viṣayaṁ} again
\end{itemize}

\begin{tcolorbox}[breakable,colback=teal!4,colframe=teal!35!black,title={Before you move on}]
What does \ml{വിഷയം} mean? (\textquotedblleft{}A subject, a topic\textquotedblright{}.) Say it once more.
\end{tcolorbox}
\section[\texorpdfstring{\ml{കണക്ക്} (kaṇakkŭ)}{കണക്ക് (kaṇakkŭ)}]{\ml{കണക്ക്} (kaṇakkŭ) — mathematics; accounts}

\label{lesson:ML-C335-kanakku}



\begin{quote}
Before the new one: say the Malayalam for a degree, then the Malayalam for a subject.
\end{quote}

\subsection*{You'll want to know: \ml{കണക്ക്}}
\textbf{\ml{കണക്ക്}} — \emph{kaṇakkŭ} — \textquotedblleft{}mathematics; accounts\textquotedblright{}.

\begin{cousinweb}
\textbf{\ml{കണക്ക്}} (\emph{kaṇakkŭ}) is of uncertain origin; DEDR, the dictionary of the Dravidian family, has no entry for it.
\end{cousinweb}

\begin{grammarlens}[title={What you've built}]
One more word for this chapter.
\end{grammarlens}

\subsection*{Your turn}
\begin{itemize}
\raggedright
  \item \emph{Say it:} \emph{kaṇakkŭ}
  \item \emph{Say it:} \emph{kaṇakkŭ}, once more
  \item \emph{Say it:} say \emph{kaṇakkŭ}
  \item \emph{Recall:} say the Malayalam for a degree, then the Malayalam for a subject, then say \emph{kaṇakkŭ} again
\end{itemize}

\begin{tcolorbox}[breakable,colback=teal!4,colframe=teal!35!black,title={Before you move on}]
What does \ml{കണക്ക്} mean? (\textquotedblleft{}Mathematics; accounts\textquotedblright{}.) Say it once more.
\end{tcolorbox}
\section[\texorpdfstring{\ml{ഗൃഹപാഠം} (gṛhapāṭhaṁ)}{ഗൃഹപാഠം (gṛhapāṭhaṁ)}]{\ml{ഗൃഹപാഠം} (gṛhapāṭhaṁ) — homework}

\label{lesson:ML-C335-grhapatham}



\begin{quote}
Before the new one: say the Malayalam for a subject, then the Malayalam for mathematics.
\end{quote}

\subsection*{You'll want to know: \ml{ഗൃഹപാഠം}}
\textbf{\ml{ഗൃഹപാഠം}} — \emph{gṛhapāṭhaṁ} — \textquotedblleft{}homework\textquotedblright{}.

\begin{cousinweb}
\textbf{\ml{ഗൃഹപാഠം}} (\emph{gṛhapāṭhaṁ}) is a modern compound of Sanskrit \emph{gṛha}, \textquotedblleft{}house\textquotedblright{} (CDIAL 4240), and \emph{pāṭha}, \textquotedblleft{}lesson\textquotedblright{} (CDIAL 8035).
\end{cousinweb}

\begin{grammarlens}[title={What you've built}]
One more word for this chapter.
\end{grammarlens}

\subsection*{Your turn}
\begin{itemize}
\raggedright
  \item \emph{Say it:} \emph{gṛhapāṭhaṁ}
  \item \emph{Say it:} \emph{gṛhapāṭhaṁ}, once more
  \item \emph{Say it:} say \emph{gṛhapāṭhaṁ}
  \item \emph{Recall:} say the Malayalam for a subject, then the Malayalam for mathematics, then say \emph{gṛhapāṭhaṁ} again
\end{itemize}

\begin{tcolorbox}[breakable,colback=teal!4,colframe=teal!35!black,title={Before you move on}]
What does \ml{ഗൃഹപാഠം} mean? (\textquotedblleft{}Homework\textquotedblright{}.) Say it once more.
\end{tcolorbox}
\section[\texorpdfstring{\ml{അക്ഷരം} (akṣaraṁ)}{അക്ഷരം (akṣaraṁ)}]{\ml{അക്ഷരം} (akṣaraṁ) — a letter (of the alphabet)}

\label{lesson:ML-C335-aksaram}



\begin{quote}
Before the new one: say the Malayalam for mathematics, then the Malayalam for homework.
\end{quote}

\subsection*{You'll want to know: \ml{അക്ഷരം}}
\textbf{\ml{അക്ഷരം}} — \emph{akṣaraṁ} — \textquotedblleft{}a letter (of the alphabet)\textquotedblright{}.

\begin{cousinweb}
\textbf{\ml{അക്ഷരം}} (\emph{akṣaraṁ}) $\leftarrow$ Sanskrit \emph{akṣara}, \textquotedblleft{}indestructible; syllable\textquotedblright{}, taken over whole as a tatsama, with Malayalam's noun ending \emph{-ṁ} (CDIAL 38). A syllable was \textquotedblleft{}the imperishable\textquotedblright{}, and became a letter.
\end{cousinweb}

\begin{grammarlens}[title={What you've built}]
One more word for this chapter.
\end{grammarlens}

\subsection*{Your turn}
\begin{itemize}
\raggedright
  \item \emph{Say it:} \emph{akṣaraṁ}
  \item \emph{Say it:} \emph{akṣaraṁ}, once more
  \item \emph{Say it:} say \emph{akṣaraṁ}
  \item \emph{Recall:} say the Malayalam for mathematics, then the Malayalam for homework, then say \emph{akṣaraṁ} again
\end{itemize}

\begin{tcolorbox}[breakable,colback=teal!4,colframe=teal!35!black,title={Before you move on}]
What does \ml{അക്ഷരം} mean? (\textquotedblleft{}A letter (of the alphabet)\textquotedblright{}.) Say it once more.
\end{tcolorbox}
